% 第 3 讲 带驯分歧的范数群
% 源：src/lecture03.txt（Raskin, 18.786 讲义，2017-08-21 版）
% 说明：原文页眉「3. NORM GROUPS WITH TAME RAMIFICATION 11/13」为 PDF 页眉残留，已删。

\section*{第 3 讲\quad 带驯分歧的范数群}

设 $K$ 是 $\operatorname{char}(K)\neq 2$ 的域。则
\[
  K^{\times}/(K^{\times})^2
  \cong\{\text{连续同态 } \operatorname{Gal}(K)\to\mathbb{Z}/2\mathbb{Z}\}
  \cong\{K \text{ 上的二次平展代数}\},
\]
这与断言 1.8 中最初的陈述对偶（这个结果是 Kummer 理论的一个雏形）。注意 $K$ 上的平展代数或是 $K\times K$，或是二次扩张 $K(\sqrt{d})/K$；前者对应 $K^{\times}/(K^{\times})^2$ 中平凡的平方陪集，后者对应由 $d\in K^{\times}$ 确定的陪集。

若 $K$ 是局部的，则 lcft 说 $\operatorname{Gal}^{\mathrm{ab}}(K)\cong\widehat{K^{\times}}$ 标准地成立。结合起来（因为这样的同态当然穿过 $\operatorname{Gal}^{\mathrm{ab}}(K)$），我们得到 $K^{\times}/(K^{\times})^2$ 有限且标准自对偶。这等价于断言存在一个「足够好」的配对
\[
  (\cdot,\cdot)\colon K^{\times}/(K^{\times})^2\times K^{\times}/(K^{\times})^2\to\{1,-1\},
\]
即它是\emph{双乘的}，满足
\[
  (a,bc)=(a,b)(a,c),\qquad (ab,c)=(a,c)(b,c),
\]
并且是\emph{非退化的}，满足条件
\[
  \text{若对一切 } b \text{ 有 } (a,b)=1,\text{ 则 } a\in(K^{\times})^2.
\]

我们曾能给出这个配对的一个简单定义，即
\[
  (a,b)=1 \iff ax^2+by^2=1 \text{ 在 } K \text{ 中有解}.
\]
注意由这个定义显然有 $(a,b)=(b,a)$，但不幸的是双乘性与非退化性都不显然，尽管本讲中我们会在许多情形证明它们成立。我们在断言 2.11 中已经证明希尔伯特符号的一个较不对称的定义成立，即对一切 $a$，
\[
  (a,b)=1 \iff b \text{ 是 } K(\sqrt{a})/K=K[t]/(t^2-a) \text{ 中的范数},
\]
而当 $a$ 是平方时，后者同构于 $K\times K$，一切元素都是范数。

在第 2 讲末尾，我们提出了下述断言，并指出重要的是这个范数子群「不能太大」（不是全部）、也「不能太小」，且 $K$ 必须是局部的。

\textbf{断言 3.1.} 这些「好性质」，即双乘性与非退化性，对希尔伯特符号成立，当且仅当对一切二次扩张 $L/K$，$N(L^{\times})\subseteq K^{\times}$ 是指数为 $2$ 的子群，即 $K^{\times}/N(L^{\times})=\mathbb{Z}/2\mathbb{Z}$。

\textit{证明.} 假设对一切二次扩张 $L/K$，$N(L^{\times})\subseteq K^{\times}$ 都是指数 $2$ 的子群。设 $a\in K^{\times}$。我们想证明
\begin{equation}\tag{3.1}
  (a,\cdot)\colon K^{\times}\to\{1,-1\}
\end{equation}
是同态，这等价于双乘性的第一个等式（另一个由对称性得到）。若 $a$ 是平方，则这很清楚，因为其像恒为 $1$（可取 $(x,y)=(1/\sqrt{a},0)$）。若 $a$ 不是平方，则令 $L:=K(\sqrt{a})$；由断言 2.11，$(a,b)=1$ 当且仅当 $b\in N(L^{\times})$。现在，与 $a$ 缩并的希尔伯特符号分解为
\[
  K^{\times}\twoheadrightarrow K^{\times}/NL^{\times}\cong\{1,-1\},
\]
其中同构是标准的，因为两个群的阶都为 $2$；我们用了范数群的指数为 $2$ 这一事实来构造保单位的最后这个 $2$ 阶群之间的双射，否则商就太大了。投影平凡地是同态。

现在证明非退化性。设 $a\notin(K^{\times})^2$。则存在某个 $b\in K^{\times}$ 不是来自 $L:=K(\sqrt{a})$ 的范数，而这当且仅当 $(a,b)=-1$，故由逆否命题，非退化性成立。

反之，注意若 $a\notin(K^{\times})^2$，则由非退化性，(3.1) 中的映射满射；又由双乘性它是同态。于是其核，即 $N(K(\sqrt{a})^{\times})\subseteq K^{\times}$，指数必为 $\#\{1,-1\}=2$。\hfill$\square$

\textbf{例 3.2.} 同样，基本情形是 $\mathbb{C}/\mathbb{R}$，其中范数群就是 $\mathbb{R}_{>0}$，指数为 $2$。

现在还需证明：

\textbf{定理 3.3.} 若 $L/K$ 是局部域的二次扩张，且 $\operatorname{char}(K)\neq 2$，则 $NL^{\times}\subseteq K^{\times}$ 是指数为 $2$ 的子群。

注意下述证明没有覆盖剩余特征为 $2$ 的分歧情形。

\textit{证明.} 令 $L:=K(\sqrt{d})$（其中 $d$ 如此前的 $a$），于是 $L$ 只依赖于 $d$ 模平方的类。则有两种情形：$v(d)=0$，当且仅当 $d\in\mathcal{O}_K^{\times}$；以及 $v(d)=1$，当且仅当 $d$ 是素元（因为可以不断约去 $\pi^2$ 的因子；注意这里 $v$ 照常是赋值）。
% [原文疑误/抽取缺字：抽出文本作 "which is true if and only if O^×_K"，缺 d∈O^×_K；已补。]

\emph{情形 1.} 此时 $d$ 是单位，且 $\sqrt{d}\in\mathcal{O}_L^{\times}$。这个扩张不一定非分歧，但我们只做非分歧情形，把分歧情形留到下周。分歧的例子是 $\mathbb{Q}_2(\sqrt{3})/\mathbb{Q}_2$（或 $\mathbb{Q}_2(\sqrt{2})$、$\mathbb{Q}_2(\sqrt{-1})$ 等）；扩张 $\mathbb{Q}_2(\sqrt{5})/\mathbb{Q}_2$ 非分歧。我们需要：
% [原文疑误：抽出文本作 "Here d is a square, and √d ∈ O^×_K"；按 v(d)=0 与下文，应为 d 是单位、√d∈O^×_L，见报告。]

\textbf{断言 3.4.} $N(\mathcal{O}_L^{\times})=\mathcal{O}_K^{\times}$；更一般地，$x\in K^{\times}$ 是范数当且仅当 $v(x)$ 为偶数（故 $K$ 中的素元不是范数）。

\textit{证明.} 我们使用「滤过技巧」。我们有下述滤过，它们被范同态保持：
\[
  \begin{CD}
    \mathcal{O}_L^{\times} @<<< 1+\mathfrak{p}_L @<<< 1+\mathfrak{p}_L^2 @<<< \cdots\\
    @V{N}VV @V{N}VV @V{N}VV\\
    \mathcal{O}_K^{\times} @<<< 1+\mathfrak{p}_K @<<< 1+\mathfrak{p}_K^2 @<<< \cdots
  \end{CD}
\]
在相伴分次项上，我们首先有
\[
  \begin{CD}
    \mathcal{O}_L^{\times}/(1+\mathfrak{p}_L) @>{N}>> \mathcal{O}_K^{\times}/(1+\mathfrak{p}_K)\\
    @V{\cong}VV @V{\cong}VV\\
    k_L^{\times} @>{N}>> k_K^{\times}.
  \end{CD}
\]
由于我们在非分歧情形，$k_L/k_K$ 与 $L/K$ 一样是二次扩张。为证明范映射在相伴分次项上满射，只需证明它在剩余域上满射，即：

\textbf{断言 3.5.} 范映射
\[
  N\colon\mathbb{F}_{q^2}^{\times}\to\mathbb{F}_q^{\times},\qquad x\mapsto x^{q+1}=\mathrm{Frob}(x)\cdot x
\]
满射（注意：因 $\operatorname{Gal}(\mathbb{F}_{q^2}/\mathbb{F}_q)\cong\mathbb{Z}/2\mathbb{Z}$，$x\mapsto x^q$ 是固定 $\mathbb{F}_q$ 的自同构；在 Galois 扩张中，域范数定义为一个元素的所有 Galois 共轭之积）。

\textit{证明 1.} 有限域的单位群必循环，故该映射对应于
\[
  \mathbb{Z}/(q^2-1)\mathbb{Z}\to\mathbb{Z}/(q-1)\mathbb{Z},\qquad n\mapsto(q+1)n.
\]
\hfill$\square$

\textit{证明 2（对 $p\neq 2$）.} 若 $x\in\mathbb{F}_q^{\times}$，则 $x=N(\sqrt{-x})$，且 $\sqrt{-x}\in\mathbb{F}_{q^2}^{\times}$，因为 $\mathbb{F}_{q^2}$ 是 $\mathbb{F}_q$ 唯一的二次扩张。\hfill$\square$

\textit{证明 3.} 我们有 $\#\mathbb{F}_{q^2}^{\times}=q^2-1$、$\#\mathbb{F}_q^{\times}=q-1$，且 $\#\ker(N)\geq(q^2-1)/(q-1)=q+1$（因像含于 $\mathbb{F}_q^{\times}$）；但多项式 $x^{q+1}-1$ 在 $\mathbb{F}_{q^2}$ 中恰有 $q+1$ 个根，因为 $\mathbb{F}_{q^2}/\mathbb{F}_q$ 可分（有限域是完美域），故 $\#\ker(N)\leq q+1$。两相结合即得 $\#\ker(N)=q+1$，从而 $N$ 满射。\hfill$\square$
% [原文疑误：抽出文本作 "#Ker(N) ≤ (q²−1)/(q−1)"，方向反了——像含于 F_q^× 只给出 ≥；已按数学含义改为 ≥，并补出两侧夹逼的结论。]
% [原文疑误：抽出文本作 "roots in F_q"，按上下文应为 F_{q²}（核在 F_{q²}^× 中），见报告。]

于是第一个相伴分次项 $\mathrm{Gr}_0$ 上的映射满射。在后续各项上，我们有
\[
  \begin{CD}
    (1+\mathfrak{p}_L^n)/(1+\mathfrak{p}_L^{n+1}) @>{N}>> (1+\mathfrak{p}_K^n)/(1+\mathfrak{p}_K^{n+1})\\
    @V{\cong}VV @V{\cong}VV\\
    k_L @>{T}>> k_K.
  \end{CD}
\]
要验证此图交换，注意由于我们假设了 $L/K$ 非分歧，$\pi$ 也是 $L$ 的素元（例如 $\mathbb{Q}_p(\sqrt{p})/\mathbb{Q}_p$ 是分歧扩张，$p$ 不再是 $\mathbb{Q}_p(\sqrt{p})$ 的素元）。于是在范映射下，
\[
  1+a\pi^n\xrightarrow{\ N\ }(1+a\pi^n)(1+\sigma a\pi^n)
  =1+(a+\sigma a)\pi^n+a\sigma a\pi^{2n}
  \equiv 1+Ta\pi^n \pmod{\pi^{n+1}}.
\]
再次地，我们提出下述断言：

\textbf{断言 3.6.} 迹映射
\[
  T\colon\mathbb{F}_{q^2}\to\mathbb{F}_q,\qquad x\mapsto x+x^q=x+\mathrm{Frob}(x)
\]
满射。

\textit{证明 1.} 迹映射的核对应 Artin--Schreier 多项式 $x^q+x$ 的根，该多项式可分，因而有 $q$ 个根，故 $\#\ker(T)=q$，且 $\#\operatorname{coker}(T)=q^2/q=q=\#\mathbb{F}_q$。\hfill$\square$

\textit{证明 2.} 若 $q$ 与 $2$ 互素，则对任意 $x\in\mathbb{F}_q$ 有 $x=T(x/2)$（如上进行）。否则若 $q=2^r$，则在 $\mathbb{F}_2$ 上 $x^2+x+1$ 是唯一的首一不可约多项式，而 $\mathbb{F}_4$ 是它的分裂域。一般地，对扩张 $\mathbb{F}_{2^{n+1}}/\mathbb{F}_{2^n}$，分裂多项式是 $x^2+x+\alpha$，其中我们选取使该多项式不可约的 $\alpha$。这恰好对一半的 $\alpha$ 选择有效，因为加法同态
\[
  \mathbb{F}_{2^n}\xrightarrow{\ x\mapsto x^2+x\ }\mathbb{F}_{2^n}
\]
的核是 $\mathbb{F}_2$，故其像的指数为 $2$。这些多项式的任何根迹都为 $1$（因为它们是首一的），而要得到任何其它迹的元素，只需乘以 $\mathbb{F}_{2^n}$ 中任意元素（因为迹映射是 $\mathbb{F}_{2^n}$-线性的）。\hfill$\square$

\textit{证明 3.} 扩张 $L/K$ 上的迹映射满射当且仅当该扩张可分，本情形如此，因为有限域是完美域。\hfill$\square$

于是，由于有限域上迹与范映射都满射，范映射在所有相伴分次项上满射，由命题 2.9 便知范映射在 $\mathcal{O}_L^{\times}$ 上满射，这证明了断言。\hfill$\square$

现在，为完成情形 1 的证明，我们想证明 $x\in K^{\times}$ 是范数当且仅当其赋值为偶数。为此注意，对任意 $y\in L^{\times}$ 有 $N(y)=y\cdot\sigma y$，且 $v(y\sigma(y))=2v(y)$，因为 $\operatorname{Gal}(L/K)$ 保持赋值。反之，只需注意 $\pi^2$ 是范数。于是若 $v(d)=0$，则 $NL^{\times}\subseteq K^{\times}$ 是指数为 $2$ 的子群，如所欲证。

\emph{情形 2.} 此时 $v(d)=1$，且仍有 $\operatorname{char}(K)\neq 2$。这保证了驯分歧，因为我们处理的是二次扩张（分歧指数不被 $p$ 整除；野分歧情形（分歧指数被 $p$ 整除）我们将在下一讲处理）。我们断言 $N(\mathcal{O}_L^{\times})\subseteq\mathcal{O}_K^{\times}$ 的指数为 $2$（具体地，$N(\mathcal{O}_L^{\times})=(\mathcal{O}_K^{\times})^2$），且存在某个 $\pi\in\mathcal{O}_K$ 既是素元又是范数。显然，这足以说明 $L^{\times}$ 的范数群在 $K^{\times}$ 中指数为 $2$。

令 $L:=K(\sqrt{d})$，其中 $v(d)=1$，故 $d$ 不是平方。则 $N(\alpha+\beta\sqrt{d})=\alpha^2-d\beta^2$，且若 $x\in(\mathcal{O}_K^{\times})^2$，则 $x=N(\sqrt{x})$，故 $x$ 是范数。反之，
\[
  v(\alpha^2-d\beta^2)=0=\min\{v(\alpha^2),v(d\beta^2)\}=v(\alpha^2),
\]
因为前者为偶数、后者为奇数，故两者不相等。由此 $\alpha,\beta\in\mathcal{O}_K^{\times}$，于是 $x=\alpha^2-d\beta^2$ 是模 $\mathfrak{p}$ 的平方，而这当且仅当 $x$ 是 $\mathcal{O}_K^{\times}$ 中的平方。最后，因为 $-d=N(\sqrt{d})$，故存在 $K$ 的一个素元是范数。
% [原文疑误：抽出文本作 "min{v(α²), v(β²+1)}"，按上下文应为 min{v(α²), v(dβ²)}，见报告。]

于是结论是：$x\in K^{\times}$ 是 $K(\sqrt{d})$ 的范数，当且仅当 $(-d)^{-v(x)}x$（一个整单位）是模 $\mathfrak{p}$ 的平方，故定理在此情形成立。\hfill$\square$

我们总结：驯分歧情形（就我们的目的而言包含非分歧情形）可以通过显式猜出 $NL^{\times}\subseteq K^{\times}$ 来处理。野分歧要棘手得多。这一切都归结为希尔伯特符号的显式公式，正如我们在习题集 1 第 2 题所见。对 $\mathbb{Q}_2$ 也有这样的公式，费点力气便可证明希尔伯特符号的一切「好」性质在此情形也成立（例如见 [Ser73]）。
